% L13a jump hidden Markov model, the shape of a jump episode next to a real crash
% (November 12, 2026, after the instructor noticed in the SPY example that "the simulated
% episodes don't look like SPY, e.g., around 2020 ... they have a diff shape"; two panels at
% his request: "the schematic on the left, and the real data plot on the right").
%
% Left, a schematic: the volatility of a simulated episode is a flat box, because every episode
% day draws its state from the same tail states; the volatility of a real crash spikes and fades.
% Right, the data, drawn like the scratch Plots.jl figure the instructor picked ("I like that
% style for the data figure"): the one-month rolling standard deviation of g (1/yr) for real SPY
% from 2020-02-20, six jump episodes of the tuned SPY model one by one, and the mean over all
% episodes, each aligned at its first day. Fig-L13a-Episode-Shape-Data.jl writes the CSV.
%
% Consistency checks (against the CSV and the script's printout):
%  - Left vertical scale: y = 0.8 + (v - 1.9) * 4.2/13, with v the rolling SD in 1/yr. Ordinary
%    days (median 1.9) sit at y = 0.8, the real peak (14.9) at y = 5.0, and the simulated
%    plateau (7.2) at y = 2.5.
%  - Left horizontal scale: 0.07 cm per trading day, with day 0 at x = 1.0. The box ends at
%    x = 8.0 (lambda = 100 days). The real curve peaks near day 30 (x = 3.1), and it is back
%    near 4.3 (y = 1.6) by day 60 (x = 5.2) and near 3 (y = 1.2) by day 100 (x = 8.0).
%  - Right panel: real peak 14.9 on day 32; simulated mean about 7 from day 20 to day 90,
%    6.3 on day 100, and 2.6 on day 150 (515 episodes from 2000 paths, the count in the
%    legend). The six examples (columns ep1 to ep6) stay between 1.1 and 10.1, inside ymax = 16.
\documentclass[tikz,border=6pt]{standalone}
% Shared style for the L13a schematics (November 12, 2026), \input by each Fig-L13a-*.tex.
% Same fonts, ink, accent colors, boxes, and arrows as the L12a schematics
% (week-12/L12a/figs/abm-diagram.tex), with the plot styles of the L10b profit diagrams
% (arrowed axes, a thick ink curve, open circle markers, dashed drops). The notebook SVG
% carries a dark palette added by theme_svg.py; every color here must be mapped there.
\usepackage[no-math]{fontspec}
\IfFontExistsTF{Arial}{\setsansfont{Arial}}{\setsansfont{TeX Gyre Heros}}
\renewcommand{\familydefault}{\sfdefault}
\usepackage{amsmath,amssymb}
\usetikzlibrary{arrows.meta,calc,positioning,decorations.pathreplacing}

\definecolor{vnink}{HTML}{202124}
\definecolor{gdblue}{HTML}{0091FF}
\definecolor{gdred}{HTML}{E02020}
\definecolor{gdgray}{HTML}{6D7278}
\definecolor{saorange}{HTML}{F08C00}
\definecolor{boxfill}{HTML}{F1F3F4}
\definecolor{orangefill}{HTML}{FFE8C7}
\definecolor{bluefill}{HTML}{E5F3FF}
\definecolor{redfill}{HTML}{FDE8E8}
\definecolor{gridgray}{HTML}{E3E5E8}

\tikzset{
  every picture/.append style={x=1cm, y=1cm, font=\sffamily\small, text=vnink},
  % boxes and arrows, as in the L12a schematics
  flow/.style={draw=vnink, line width=0.9pt, -{Latex[length=2.4mm, width=1.7mm]}},
  feedback/.style={draw=saorange, line width=1.2pt, -{Latex[length=2.6mm, width=1.9mm]}, rounded corners=6pt},
  sumbox/.style={draw=vnink, fill=boxfill, line width=0.8pt, rounded corners=3pt, inner sep=5pt, align=center},
  note/.style={text=gdgray, font=\sffamily\footnotesize, align=center},
  state/.style={circle, draw=vnink, fill=white, line width=1.0pt, minimum size=1.1cm, inner sep=0pt, font=\sffamily\normalsize},
  % plots, as in the L10b profit diagrams
  axis/.style={draw=vnink, line width=0.8pt, -{Latex[length=2.6mm, width=1.8mm]}},
  spine/.style={draw=vnink, line width=0.8pt},
  curve/.style={draw=vnink, line width=1.6pt, line join=round, line cap=round},
  series/.style={draw=gdgray, line width=0.9pt, line join=round},
  drop/.style={draw=gdgray, line width=0.8pt, dash pattern=on 3pt off 2pt},
  pt/.style={circle, draw=vnink, fill=white, line width=1.0pt, inner sep=0pt, minimum size=2.0mm},
}

\usepackage{pgfplots}
\pgfplotsset{compat=1.18}

\begin{document}
\begin{tikzpicture}

  % ================= left: the schematic =================
  \node[font=\sffamily\bfseries] at (5.5, 6.25) {Schematic};

  % axes and the level of ordinary days
  \draw[axis] (-0.15, 0) -- (11.2, 0);
  \draw[axis] (0, -0.15) -- (0, 5.8);
  \node[anchor=north] at (5.5, -0.12) {trading days since the start};
  \node[rotate=90, anchor=south] at (-0.2, 2.8) {volatility};
  \draw[drop] (0, 0.8) -- (10.9, 0.8);
  \node[note, anchor=north east] at (10.9, 0.76) {ordinary days};

  % simulated episode: a flat box for about lambda days
  \fill[redfill] (1.0, 0.8) rectangle (8.0, 2.5);
  \draw[draw=gdred, line width=1.6pt, line join=round] (1.0, 0.8) -- (1.0, 2.5) -- (8.0, 2.5) -- (8.0, 0.8);
  \node[text=gdred, align=center, anchor=south] at (6.4, 2.55) {simulated episode:\\flat for about $\lambda$ days};
  \draw[draw=gdred, line width=0.8pt, {Latex[length=2.0mm, width=1.4mm]}-{Latex[length=2.0mm, width=1.4mm]}] (1.0, 0.4) -- (8.0, 0.4);
  \node[text=gdred, fill=white, inner sep=2pt, font=\sffamily\footnotesize] at (4.5, 0.4) {about $\lambda$ days};

  % real crash: spikes, then fades
  \draw[curve] plot[smooth, tension=0.6] coordinates {
    (1.0, 0.8) (1.6, 1.4) (2.1, 3.0) (2.5, 4.4) (3.1, 5.0) (3.5, 4.5) (3.9, 3.1)
    (4.4, 2.1) (5.2, 1.6) (6.2, 1.4) (8.0, 1.2) (9.5, 1.0) (10.6, 0.9)};
  \node[align=left, anchor=west] at (3.55, 4.75) {real crash (SPY, 2020):\\spikes, then fades};

  % ================= right: the data =================
  \node[font=\sffamily\bfseries] at (18.35, 6.25) {SPY in 2020 and the tuned model};
  \begin{axis}[
      at={(13.6cm, 0cm)}, anchor=south west, scale only axis, width=9.5cm, height=5.8cm,
      set layers, % without layers, the x=1cm and y=1cm picture style of jumphmm-diagram.tex stretches the frame and the grid to data units
      xmin=-35, xmax=205, ymin=0, ymax=16,
      xtick={0, 50, 100, 150, 200}, ytick={0, 4, 8, 12, 16},
      axis background/.style={fill=boxfill},
      axis line style={draw=vnink, line width=0.8pt},
      grid=major, grid style={draw=gridgray, line width=0.5pt},
      tick style={draw=vnink}, tick align=inside,
      every tick label/.append style={font=\sffamily\footnotesize, text=vnink},
      /pgf/number format/assume math mode=true,
      xlabel={trading days since the start}, ylabel={one-month rolling SD of $g$ (1/yr)},
      every axis label/.append style={font=\sffamily\small, text=vnink},
      legend pos=north east, legend style={draw=none, fill=none, font=\sffamily\footnotesize, text=vnink},
      legend cell align=left,
    ]
    % six simulated episodes, one by one (the first one carries the legend entry)
    \addplot[draw=gdred, line width=0.6pt, draw opacity=0.5] table[x=day, y=ep1, col sep=comma] {Fig-L13a-Episode-Shape-Data.csv};
    \addlegendentry{simulated episodes (six examples)}
    \addplot[draw=gdred, line width=0.6pt, draw opacity=0.5, forget plot] table[x=day, y=ep2, col sep=comma] {Fig-L13a-Episode-Shape-Data.csv};
    \addplot[draw=gdred, line width=0.6pt, draw opacity=0.5, forget plot] table[x=day, y=ep3, col sep=comma] {Fig-L13a-Episode-Shape-Data.csv};
    \addplot[draw=gdred, line width=0.6pt, draw opacity=0.5, forget plot] table[x=day, y=ep4, col sep=comma] {Fig-L13a-Episode-Shape-Data.csv};
    \addplot[draw=gdred, line width=0.6pt, draw opacity=0.5, forget plot] table[x=day, y=ep5, col sep=comma] {Fig-L13a-Episode-Shape-Data.csv};
    \addplot[draw=gdred, line width=0.6pt, draw opacity=0.5, forget plot] table[x=day, y=ep6, col sep=comma] {Fig-L13a-Episode-Shape-Data.csv};
    % the mean over all episodes, and real SPY on top
    \addplot[draw=gdred, line width=1.8pt] table[x=day, y=sim_mean, col sep=comma] {Fig-L13a-Episode-Shape-Data.csv};
    \addlegendentry{simulated episodes (mean of 515)}
    \addplot[draw=vnink, line width=1.8pt] table[x=day, y=real, col sep=comma] {Fig-L13a-Episode-Shape-Data.csv};
    \addlegendentry{real SPY, from February 20, 2020}
    % the start of the episode
    \draw[draw=vnink, line width=0.6pt, dash pattern=on 4pt off 3pt] (axis cs:0, 0) -- (axis cs:0, 16);
  \end{axis}

\end{tikzpicture}
\end{document}
